% pw1-plot-maze TEXT-INSERTS.tex (PW-1 items 16.1, 1.5, 8). Not installed; a later spoke places these.

% ---------------------------------------------------------------- (a) the panel figures, read once (item 16.1)
% Place in the paragraph of Sec. sec:maze-wm that introduces Figs. fig:cmazet-families -- fig:cmazet-door-tint.
% Replaces the paragraph that was printed under every panel image (cmazet_panels.KEY_LINE, no longer drawn).
% SRC: w21_figures/cmazet_panels.py (draw_state, fact_strip, facts_shown); chlu/exp_w15_cmaze_t.py (outcome code).
\paragraph{Reading the CMAZE-T panel figures.} Each panel is a $5\times5$ crop of the maze centred on
the agent's cell at $s_t$. Walls are dark, tinted cells pale in their colour, doors framed in their
colour with a bar across while closed, untaken keys are diamonds, the goal a cross, and the agent a
black disc whose centre shows the colour of the key in its pocket (white when the pocket is empty).
The first column is $s_t$ with the action (an arrow, or a ring for pick-up), the second the engine's
$s_{t+1}$, and each further column the $s_{t+1}$ that a model's emitted outcome produces. The engine's
judge does not score the picture: it scores four emitted facts, legality, the move, the key event and
the pocket after the move, and a picture need not show all four. So a strip under every $s_{t+1}$
prints the emitted facts (under the engine's column, the true ones). Its first line reads legal or
illegal. Its second line shows the move ($\uparrow$ $\rightarrow$ $\downarrow$ $\leftarrow$, or a dot
for staying), the key event ($+$ and a colour for a key picked up, $-$ and a colour for a key used at a
door, blank for none) and the pocket after the move (the disc; white is empty). An illegal emission
leaves the agent in place and is scored on its pocket alone, so only the pocket is shown. The tick or
cross on a panel is the judge's verdict on all four facts, so two panels that show the same picture
and the same strip never carry opposite marks.

% the one sentence each panel caption uses to point here (already in CAPTIONS.tex):
%   The strip under each $s_{t+1}$ is read as described in Sec.~\ref{sec:maze-wm}.

% ---------------------------------------------------------------- (b) the hold-out mechanism (items 1.5, 8)
% Replaces the task-section paragraph on hard seeds and the per-seed ceiling (DRAFT-v2.tex 1138-1150,
% "The measured per-seed ceiling, ..."). No seed is labelled; no ceiling number.
% SRC: w21-maze-t-seed4/FINDINGS.md (mechanism: same-colour pocket/tint rows removed by the draw; identity vs
% relation); chlu/exp_w15_cmaze_t.py (RULE B: pocket colour equal to tint colour makes the move illegal).
\paragraph{Why a hold-out draw can remove a whole colour relation.} The hold-out removes word
combinations from training at random. On some draws, for some colour, it removes every training
transition in which the pocket and the target tint share that colour, the combination that decides
whether a move is legal. A model that codes colour by identity then has no training evidence for the
held-out transitions of that colour and cannot answer them, however it is trained. A model that is
also given the relation ``pocket colour equals tint colour'' can: the relation is learned from the
other colours, where it is seen, and applies unchanged to the colour whose rows were removed. This is
what the relational columns add (Fig.~\ref{fig:maze-wm}(c)), and it is why every seed is reported
separately.
